%!TEX program = xelatex
\documentclass[11pt, a4paper]{article}

% AA LINGUA E GESTIONE FONT (Specifica per XeLaTeX)
\usepackage{fontspec}
\setmainfont{Latin Modern Roman}
\usepackage[italian]{babel}

% AA LAYOUT DI PAGINA E TIPOGRAFIA (Ottimizzato per Rilegatura)
\usepackage[left=2.4cm, right=1.6cm, top=2cm, bottom=2cm]{geometry}

% AA Spaziatura tra paragrafi controllata (compatta ed evita salti enormi)
\usepackage[skip=3pt plus 1pt minus 1pt]{parskip}
\setlength{\parindent}{0pt}

% AA Microtype per XeLaTeX
\usepackage[protrusion=true]{microtype}
\raggedbottom

% AA GESTIONE DEI COLORI
\usepackage{xcolor}

% AA MATEMATICA E SIMBOLI
\usepackage{amsmath}
\usepackage{mathtools}
\usepackage{silence}
\usepackage[warnings-off={mathtools-overbracket,mathtools-colon}]{unicode-math}
\setmathfont{Latin Modern Math}

% BB OTTIMIZZAZIONE MATRICI LUNGHE PER LA STAMPA
\setlength{\arraycolsep}{10pt}

% AA Spaziature verticali attorno alle formule matematiche
\setlength{\abovedisplayskip}{3pt plus 1pt minus 1pt}
\setlength{\belowdisplayskip}{3pt plus 1pt minus 1pt}
\setlength{\abovedisplayshortskip}{0pt plus 1pt}
\setlength{\belowdisplayshortskip}{2pt plus 1pt minus 1pt}

% AA PERSONALIZZAZIONE TITOLI
\usepackage{titlesec}
\titleformat{\section}{\normalfont\Large\bfseries\color{blue!80!black}}{\thesection}{1em}{}
\titleformat{name=\section,numberless}{\normalfont\Large\bfseries\color{blue!80!black}}{}{0pt}{}

\titleformat{\subsection}{\normalfont\large\bfseries}{\thesubsection}{1em}{}
\titleformat{name=\subsection,numberless}{\normalfont\large\bfseries}{}{0pt}{}

\titleformat{\subsubsection}{\normalfont\normalsize\bfseries}{\thesubsubsection}{1em}{}
\titleformat{name=\subsubsection,numberless}{\normalfont\normalsize\bfseries}{}{0pt}{}

\titlespacing*{\section}{0pt}{1.2ex plus 0.4ex minus 0.2ex}{0.6ex plus 0.2ex}
\titlespacing*{\subsection}{0pt}{1ex plus 0.3ex minus 0.2ex}{0.5ex plus 0.2ex}
\titlespacing*{\subsubsection}{0pt}{0.8ex plus 0.2ex minus 0.1ex}{0.4ex plus 0.1ex}

\titleformat{\paragraph}[block]{\normalfont\normalsize\bfseries\color{black!80!black}\raggedright}{}{0pt}{}
\titlespacing*{\paragraph}{0pt}{1.2ex plus 0.3ex minus 0.2ex}{0.4ex}

% AA INTEGRAZIONE CODICE MATLAB
\usepackage{matlab}

% BB Riconfigurazione Automatica Titoli Esportati da MATLAB
\renewcommand{\matlabheading}[1]{\section*{#1}}
\renewcommand{\matlabheadingtwo}[1]{\subsection*{#1}}
\renewcommand{\matlabheadingthree}[1]{\subsubsection*{#1}}

% BB Formule E Matrici Simboliche
\renewenvironment{matlabsymbolicoutput}{%
  \par\vspace{2pt}%
  \begingroup
  \centering
}{%
  \par\vspace{2pt}%
  \endgroup
}

% BB Output Testuale Console
\renewenvironment{matlaboutput}{%
  \par\vspace{1pt}%
  \Verbatim[xleftmargin=10pt, formatcom=\color{black!75}, fontsize=\footnotesize]%
}{%
  \endVerbatim
  \vspace{1pt}%
}

% BB Box Codice (Stile Moderno con Barra Laterale Colorata)
\lstset{
  backgroundcolor={},
  aboveskip=0pt,
  belowskip=0pt,
  basicstyle=\ttfamily\footnotesize,
  xleftmargin=0pt
}

\mdfdefinestyle{matlabcode}{%
  topline=false,
  bottomline=false,
  rightline=false,
  leftline=true,
  linewidth=2.5pt,
  linecolor=blue!75!black,
  backgroundcolor=black!3,
  innertopmargin=4pt,
  innerbottommargin=4pt,
  innerleftmargin=9pt,
  innerrightmargin=6pt,
  skipabove=3pt,
  skipbelow=3pt
}

\allowdisplaybreaks

% AA LINK, METADATI ED ELENCHI
\usepackage{hyperref}
\hypersetup{
    colorlinks=true,
    linkcolor=blue!80!black,
    urlcolor=blue!80!black,
    citecolor=blue!80!black,
    pdftitle={Calcolo Jacobiano Geometrico},
    pdfauthor={Francesco Lo Verde},
}
\usepackage{bookmark}

% BB Spaziatura Elenchi (Arieggiata per la Stampa)
\usepackage{enumitem}
\setlist{itemsep=3.5pt plus 1pt minus 1pt, topsep=4pt plus 1pt minus 1pt, partopsep=0pt, parsep=1pt}


\author{Francesco Lo Verde}
\begin{document}

\label{T_82ff}
\matlabtitle{Calcolo E Analisi Dello Jacobiano Geometrico}

\begin{par}
\begin{flushleft}
Il calcolo dello Jacobiano geometrico $J(q)$, uno strumento fondamentale nell'analisi cinematica dei manipolatori seriali. Lo \textbf{Jacobiano} lega le velocità dei giunti $\dot{q}$ alle velocità cartesiane dell'end-effector (lineare $v_E$ e angolare $\omega_E$) attraverso la relazione:
\end{flushleft}
\end{par}

\label{H_91e8}
\matlabheadingthree{$\left(\begin{array}{c}
v_E \\
\omega_E
\end{array}\right)=J(q)\dot{q}$}

\begin{par}
\begin{flushleft}
Il metodo sistematico si basa sulla scomposizione dello Jacobiano in colonne $J_i$, dove ciascuna colonna descrive l'influenza del giunto i-esimo sulla velocità dell'end-effector. Per calcolare lo Jacobiano geometrico, è necessario conoscere:
\end{flushleft}
\end{par}

\begin{itemize}
\setlength{\itemsep}{-1ex}
   \item{\begin{flushleft} La \textbf{geometria del manipolatore}, descritta dai parametri di \textit{Denavit-Hartenberg (DH)}. \end{flushleft}}
   \item{\begin{flushleft} Il \textbf{tipo di giunto} (rotoidale o prismatico). \end{flushleft}}
   \item{\begin{flushleft} Il \textbf{numero di giunti }del manipolatore \end{flushleft}}
\end{itemize}

\begin{par}
\begin{flushleft}
Si definiscono i parametri del manipolatore RR planare. Le lunghezze dei bracci ($l_1 ,l_2$) e le variabili di giunto ($q_1 ,q_2$) sono dichiarate come variabili simboliche per consentire un calcolo analitico.
\end{flushleft}
\end{par}

\begin{matlabcode}
clc; clear; close all;
% Variabili Di Giunto
syms q1 q2 real;

% Lunghezza Bracci
syms a1 a2 real;
\end{matlabcode}

\begin{par}
\begin{flushleft}
La cinematica del manipolatore si ottiene dalla tabella dei parametri di Denavit-Hartenberg:
\end{flushleft}
\end{par}

\begin{matlabcode}
% Definizione Tabelle DH Struttura
DH_Table = [q1, 0, 0, a1;
            q2, 0, 0, a2];
% q1 = pi/4;
% q2 = pi/2;
% a1 = 1;
% a2 = 0.8
\end{matlabcode}

\begin{par}
\begin{flushleft}
Per calcolare lo Jacobiano, è necessario conoscere la posizione e l'orientamento di ogni terna rispetto alla terna base ($RF_0$). Queste informazioni sono contenute nelle matrici di trasformazione omogenea $^{i-1} T_i$, che lega la terna $i$ alla terna $i-1$, che è calcolata a partire dalla riga i-esima della tabella DH. Per convenzione la matrince $^1 T_1 =I$.
\end{flushleft}
\end{par}

\begin{matlabcode}
% Calcolo Matrici Trasformazione E Rotazione Struttura
[T, R] = Compute_DH_Transformations(DH_Table);
\end{matlabcode}

\begin{par}
\begin{flushleft}
La matrice di trasformazione cumulativa $^0 T_2$, descrive la posa (posizione e orientamento) della terna $2$ rispetto alla terna base $0$.
\end{flushleft}
\end{par}

\begin{matlabcode}
T02 = simplify(T(:,:,1) * T(:,:,2));
fprintf('Matrice Di Trasformazione Omogenea ^{0}T_(2)');
disp(T02);
\end{matlabcode}
\begin{matlaboutput}
Matrice Di Trasformazione Omogenea ^{0}T_(2)
\end{matlaboutput}
\begin{matlabsymbolicoutput}
\hskip1em $\displaystyle \left(\begin{array}{cccc}
\cos \left(q_1 +q_2 \right) & -\sin \left(q_1 +q_2 \right) & 0 & a_2 \,\cos \left(q_1 +q_2 \right)+a_1 \,\cos \left(q_1 \right)\\
\sin \left(q_1 +q_2 \right) & \cos \left(q_1 +q_2 \right) & 0 & a_2 \,\sin \left(q_1 +q_2 \right)+a_1 \,\sin \left(q_1 \right)\\
0 & 0 & 1 & 0\\
0 & 0 & 0 & 1
\end{array}\right)$
\end{matlabsymbolicoutput}
\newpage
\begin{par}
\begin{flushleft}
Lo Jacobiano geometrico $J$ è una matrice 6xn le cui colonne $J_i$ sono calcolate secondo le specifiche formule, a seconda del tipo di giunto:
\end{flushleft}
\end{par}

\begin{par}

\begin{table}[h]
    \centering
    \renewcommand{\arraystretch}{1.2} % aumenta leggermente lo spazio tra le righe
    \begin{tabular}{|l|c|c|}
        \hline
        \textbf{Tipo di Giunto} & \textbf{Contributo Lineare} \(\mathbf{J}_{L_{i}}\) & \textbf{Contributo Angolare} \(\mathbf{J}_{A_{i}}\) \\
        \hline
        \textit{Prismatico} & \(z_{i-1}\) & \(\mathbf{0}_{3 \times 1}\) \\
        \hline
        \textit{Rotoidale} & \(z_{i-1} \times (p_{E} - p_{i-1})\) & \(z_{i-1}\) \\
        \hline
    \end{tabular}
\end{table}

\end{par}

\begin{itemize}
   \item{\begin{flushleft} $z_{i-1}$ è il versore dell'asse di rotazione del giunto $i$ espresso nelle coordinate della terna base. \end{flushleft}}
   \item{\begin{flushleft} $p_{i-1}$ è il vettore posizione dell'origine della terna $i-1$, espresso rispetto alla terna base. \end{flushleft}}
   \item{\begin{flushleft} $p_E$ è il vettore posizione dell'end-effector, espresso rispetto alla terna base. \end{flushleft}}
\end{itemize}
\vspace{5mm}
\begin{matlabcode}
% Estrazioni Matrici Singole
T0 = eye(4);
T1 = T(:,:,1);
T2 = T02;
\end{matlabcode}

\begin{par}
\begin{flushleft}
Il vettore $z_{i-1}$ corrisponde alla \textbf{terza colonna} della matrice di rotazione 3x3 contenuta in $T_{i-1}$.
\end{flushleft}
\end{par}

\begin{par}
\begin{flushleft}
Il vettore $p_{i-1}$ corrisponde alla \textbf{quarta colonna} (il vettore di traslazione) della matrice $T_{i-1}$.
\end{flushleft}
\end{par}

\begin{matlabcode}
% Valori Per J1
z0 = T0(1:3,3);
fprintf('Vettore z0:');
disp(z0);
\end{matlabcode}
\begin{matlaboutput}
Vettore z0:
\end{matlaboutput}
\begin{matlaboutput}
     0
     0
     1
\end{matlaboutput}
\begin{matlabcode}

p0 = T0(1:3,4);
fprintf('Vettore p0:');
disp(p0);
\end{matlabcode}
\begin{matlaboutput}
Vettore p0:
\end{matlaboutput}
\begin{matlaboutput}
     0
     0
     0
\end{matlaboutput}
\begin{matlabcode}

% Valori Per J2
z1 = T1(1:3,3);
fprintf('Vettore z1');
disp(z1);
\end{matlabcode}
\begin{matlaboutput}
Vettore z1
\end{matlaboutput}
\begin{matlabsymbolicoutput}
\hskip1em $\displaystyle \left(\begin{array}{c}
0\\
0\\
1
\end{array}\right)$
\end{matlabsymbolicoutput}
\begin{matlabcode}

p1 = T1(1:3,4);
fprintf('Vettore p1:');
disp(p1);
\end{matlabcode}
\begin{matlaboutput}
Vettore p1:
\end{matlaboutput}
\begin{matlabsymbolicoutput}
\hskip1em $\displaystyle \left(\begin{array}{c}
a_1 \,\cos \left(q_1 \right)\\
a_1 \,\sin \left(q_1 \right)\\
0
\end{array}\right)$
\end{matlabsymbolicoutput}
\begin{matlabcode}
% Valori Per JE
zE = T2(1:3,3);
fprintf('Vettore z2:');
disp(zE);
\end{matlabcode}
\begin{matlaboutput}
Vettore z2:
\end{matlaboutput}
\begin{matlabsymbolicoutput}
\hskip1em $\displaystyle \left(\begin{array}{c}
0\\
0\\
1
\end{array}\right)$
\end{matlabsymbolicoutput}
\begin{matlabcode}
pE = T2(1:3,4);
fprintf('Vettore p2:');
disp(pE);
\end{matlabcode}
\begin{matlaboutput}
Vettore p2:
\end{matlaboutput}
\begin{matlabsymbolicoutput}
\hskip1em $\displaystyle \left(\begin{array}{c}
a_2 \,\cos \left(q_1 +q_2 \right)+a_1 \,\cos \left(q_1 \right)\\
a_2 \,\sin \left(q_1 +q_2 \right)+a_1 \,\sin \left(q_1 \right)\\
0
\end{array}\right)$
\end{matlabsymbolicoutput}
\begin{matlabcode}
% Contributi Angolari
J_A1 = z0;
J_A2 = z1;

% Contributi Lineari
J_L1 = cross(z0, pE - p0);
J_L2 = cross(z1, pE - p1);

% Colonne
J1 = [J_L1; J_A1];
J2 = [J_L2; J_A2];

fprintf('Vettore J1:');
disp(J1);
\end{matlabcode}
\begin{matlaboutput}
Vettore J1:
\end{matlaboutput}
\begin{matlabsymbolicoutput}
\hskip1em $\displaystyle \left(\begin{array}{c}
-a_2 \,\sin \left(q_1 +q_2 \right)-a_1 \,\sin \left(q_1 \right)\\
a_2 \,\cos \left(q_1 +q_2 \right)+a_1 \,\cos \left(q_1 \right)\\
0\\
0\\
0\\
1
\end{array}\right)$
\end{matlabsymbolicoutput}
\begin{matlabcode}
fprintf('Vettore J2:');
disp(J2);
\end{matlabcode}
\begin{matlaboutput}
Vettore J2:
\end{matlaboutput}
\begin{matlabsymbolicoutput}
\hskip1em $\displaystyle \left(\begin{array}{c}
-a_2 \,\sin \left(q_1 +q_2 \right)\\
a_2 \,\cos \left(q_1 +q_2 \right)\\
0\\
0\\
0\\
1
\end{array}\right)$
\end{matlabsymbolicoutput}

\newpage
\begin{par}
\begin{flushleft}
A questo punto concatenando i due vettori si ottiene lo \textbf{Jacobiano Geometrico:}
\end{flushleft}
\end{par}

\begin{matlabcode}
% Costruzione Jacobiano Geometrico
J_q = [J1, J2];
J_q = simplify(J_q);

fprintf('Jacobiano Geometrico Manipolatore RR Planare Simbolico');
disp(J_q);
\end{matlabcode}
\begin{matlaboutput}
Jacobiano Geometrico Manipolatore RR Planare Simbolico
\end{matlaboutput}
\begin{matlabsymbolicoutput}
\hskip1em $\displaystyle \left(\begin{array}{cc}
-a_2 \,\sin \left(q_1 +q_2 \right)-a_1 \,\sin \left(q_1 \right) & -a_2 \,\sin \left(q_1 +q_2 \right)\\
a_2 \,\cos \left(q_1 +q_2 \right)+a_1 \,\cos \left(q_1 \right) & a_2 \,\cos \left(q_1 +q_2 \right)\\
0 & 0\\
0 & 0\\
0 & 0\\
1 & 1
\end{array}\right)$
\end{matlabsymbolicoutput}

\vspace{1em}

\matlabheadingtwo{Analisi Singolarità}

\begin{par}
\begin{flushleft}
Le singolarità cinematiche si verificano nelle configurazioni in cui il manipolatore perde la capacità di muoversi in una o più direzioni dello spazio operativo. Matematicamente, questo accade quando la matrice Jacobiana perde di rango. Per un manipolatore planare, si è interessati principalmente al movimento nel piano XY, quindi si analizza la sottomatrice dello Jacobiano che lega le velocità dei giunti alle velocità lineari dell'end-effector.
\end{flushleft}
\end{par}

\begin{matlabcode}
% Estrazione Jacobiano Lineare (Prime Due Righe)
J_L = J_q(1:2,1:2);
fprintf('Jacobiano Lineare Manipolatore RR Planare Simbolico:');
disp(J_L);
\end{matlabcode}
\begin{matlaboutput}
Jacobiano Lineare Manipolatore RR Planare Simbolico:
\end{matlaboutput}
\begin{matlabsymbolicoutput}
\hskip1em $\displaystyle \left(\begin{array}{cc}
-a_2 \,\sin \left(q_1 +q_2 \right)-a_1 \,\sin \left(q_1 \right) & -a_2 \,\sin \left(q_1 +q_2 \right)\\
a_2 \,\cos \left(q_1 +q_2 \right)+a_1 \,\cos \left(q_1 \right) & a_2 \,\cos \left(q_1 +q_2 \right)
\end{array}\right)$
\end{matlabsymbolicoutput}

\begin{par}
\begin{flushleft}
Le configurazioni singolari si trovano cercando i valori delle variabili di giunto $q$ che annullano il determinante della Jacobiana Lineare.
\end{flushleft}
\end{par}

\begin{matlabcode}
% Calcolo Determinante Jacobiano Lineare
detJL = det(J_L);
detJL = simplify(detJL);

fprintf('Determinante Jacobiano Lineare:');
disp(detJL);
\end{matlabcode}

\begin{matlaboutput}
Determinante Jacobiano Lineare:
\end{matlaboutput}
\begin{matlabsymbolicoutput}
\hskip1em $\displaystyle a_1 \,a_2 \,\sin \left(q_2 \right)$
\end{matlabsymbolicoutput}

\begingroup
\mdfsetup{leftmargin=-5mm, rightmargin=-5mm, innerrightmargin=4pt}
\begin{matlabcode}
\begin{matlabcode}
[solDetJL, params, conditions] = solve(detJL == 0, q2, 'ReturnConditions', true);
% Valori pari di k corrispondono al robot completamente esteso (q2 = 0).
% Valori dispari di k corrispondono al robot completamente ripiegato (q2 = $\pi$).
% Basta considerare k = 0 e k = 1
% params = k
solDetJL_1 = subs(solDetJL, params, 0);
solDetJL_2 = subs(solDetJL, params, 1);
\end{matlabcode}
\endgroup

\newpage
\begin{matlabcode}
fprintf('Valori Per Cui Determinante Jacobiano Nullo:');
disp(solDetJL_1);
disp(solDetJL_2);
\end{matlabcode}
\begin{matlaboutput}
Valori Per Cui Determinante Jacobiano Nullo:
\end{matlaboutput}
\begin{matlabsymbolicoutput}
\hskip1em $\displaystyle 0$
\hskip1em $\displaystyle \pi $
\end{matlabsymbolicoutput}

\begin{par}
\begin{flushleft}
Il risultato dell'equazione $det(J_L )=0$ fornisce le configurazioni di singolarità del manipolatore. Per il manipolatore RR planare, queste condizioni corrispondono fisicamente alle configurazioni in cui i due bracci sono allineati, ovvero quando il secondo giunto ha un angolo di 0 o  radianti. In queste pose, l'end-effector non può muoversi istantaneamente in direzione radiale.
\end{flushleft}
\end{par}

\matlabheadingtwo{Calcolo Numerico}

\begin{matlabcode}
%% Calcolo Numerico Configurazione Specifica

% Definizione Valori Numerici
q1_val = pi/4;
q2_val = pi/2;
a1_val = 1;
a2_val = 0.8;

% Sostituzione
J_q_num = subs(J_q, {q1, q2, a1, a2}, {q1_val, q2_val, a1_val, a2_val});
J_L_num = subs(J_L, {q1, q2, a1, a2}, {q1_val, q2_val, a1_val, a2_val});
detJL_num = subs(detJL, {q1, q2, a1, a2}, {q1_val, q2_val, a1_val, a2_val});

fprintf('Jacobiano Geometrico Numerico J_q:');
disp(double(J_q_num));
\end{matlabcode}
\begin{matlaboutput}
Jacobiano Geometrico Numerico J_q:
\end{matlaboutput}
\begin{matlaboutput}
   -1.2728   -0.5657
    0.1414   -0.5657
         0         0
         0         0
         0         0
    1.0000    1.0000
\end{matlaboutput}
\begin{matlabcode}
fprintf('Jacobiano Lineare Numerico J_L:');
disp(double(J_L_num));
\end{matlabcode}
\begin{matlaboutput}
Jacobiano Lineare Numerico J_L:
   -1.2728   -0.5657
    0.1414   -0.5657
\end{matlaboutput}
\begin{matlabcode}
fprintf('Valore del Determinante dello Jacobiano Lineare:');
disp(double(detJL_num));
\end{matlabcode}
\begin{matlaboutput}
Valore del Determinante dello Jacobiano Lineare:
    0.8000
\end{matlaboutput}
\begin{matlabcode}
% Verifica Singolarità
if abs(double(detJL_num)) < 1e-10
    fprintf('Configurazione Singolare');
else
    fprintf('Configurazione NON Singolare.\n');
end
\end{matlabcode}
\begin{matlaboutput}
Configurazione NON Singolare.
\end{matlaboutput}

\end{document}
